\documentclass[10pt, a4paper]{article}

% Packages
\usepackage[utf8]{inputenc}
\usepackage[T1]{fontenc}
\usepackage[margin=2.5cm]{geometry}
\usepackage{xcolor}
\usepackage{titlesec}
\usepackage{enumitem}
\usepackage{hyperref}
\usepackage{fontawesome5}
\usepackage{charter}
\usepackage{parskip}

% Primary Color Definition: RGB(0, 51, 102)
\definecolor{primarycolor}{RGB}{0, 51, 102}

% Hyperlink Setup
\hypersetup{
    colorlinks=true,
    linkcolor=primarycolor,
    urlcolor=primarycolor,
}

\begin{document}

% --- SENDER INFO ---
\begin{flushright}
    \textbf{Kartavya Niraj Dikshit} \\
    Halbmondstraße 7, 91054 Erlangen \\
    \href{mailto:kartavya.dikshit@gmail.com}{kartavya.dikshit@gmail.com} \\
    (+49) 15510940829 \\
    \small \href{https://www.linkedin.com/in/kartavya-dikshit}{linkedin.com/in/kartavya-dikshit} \\
    \href{https://github.com/KartavyaDikshit}{github.com/KartavyaDikshit}
\end{flushright}

\vspace{10pt}

% --- RECIPIENT INFO ---
Siemens AG \\
Advanta – Internal Services \\
Recruiting Team \\
Otto-Hahn-Ring 6 \\
81739 München

\vspace{20pt}

% --- DATE & SUBJECT ---
Erlangen, 17. März 2026

\vspace{10pt}
\textbf{\large Bewerbung als Werkstudent (w/m/d) Advanced Analytics and Artificial Intelligence} \\
\textbf{Job ID: 497378}

\vspace{15pt}

% --- SALUTATION ---
Sehr geehrte Damen und Herren,

mit großem Interesse verfolge ich die Innovationskraft von Siemens Advanta bei der Transformation industrieller Daten in wertschöpfende KI-Lösungen. Als Master-Student der Data Science an der Friedrich-Alexander-Universität Erlangen-Nürnberg (FAU) mit den Schwerpunkten Scalable Data Architectures und Advanced Machine Learning suche ich eine Herausforderung, bei der ich meine akademische Spezialisierung direkt in einem weltweit führenden Technologieunternehmen einbringen kann.

Mein Studium an der FAU stattet mich mit dem theoretischen Fundament aus, das für die von Ihnen ausgeschriebene Position essentiell ist. Besonders in den Modulen zu skalierbaren Datenarchitekturen und Deep Learning habe ich gelernt, komplexe Algorithmen auf große Datenmengen anzuwenden – eine Fähigkeit, die ich bereits in einem Projekt zur prädiktiven Wartung mittels Azure und Snowflake erfolgreich unter Beweis gestellt habe. Dabei konnte ich die Ausfallzeiten industrieller Anlagen durch optimierte ETL-Prozesse und präzise Telemetrie-Analysen signifikant reduzieren.

Durch meine bisherigen Praktika und freiberuflichen Tätigkeiten verfüge ich über fundierte Kenntnisse in Python (Pandas, Scikit-learn) sowie SQL. Ich bin erfahren darin, Datenpipelines in Cloud-Umgebungen wie Azure und AWS zu entwickeln und diese durch interaktive Dashboards in Tableau oder Power BI für Stakeholder greifbar zu machen. Mein Ziel ist es nicht nur, Modelle zu entwickeln, sondern "Actionable Insights" zu generieren, die Geschäftsprozesse nachhaltig verbessern.

Siemens Advanta reizt mich besonders durch die einzigartige Kombination aus technischer Tiefe und operativem Einfluss. Ich schätze die agile Zusammenarbeit in cross-funktionalen Teams und bringe durch die Leitung technischer Workshops die nötige Kommunikationsstärke mit, um komplexe Sachverhalte verständlich zu vermitteln. Meine fließenden Englischkenntnisse (C2) sowie meine fortgeschrittenen Deutschkenntnisse (B1, mit dem Ziel der professionellen Sprachkompetenz) ermöglichen mir eine reibungslose Integration in Ihr internationales Team.

Gerne möchte ich Sie in einem persönlichen Gespräch von meiner Motivation und meinen technischen Fähigkeiten überzeugen. Ich stehe Ihnen für 15 bis 20 Stunden pro Woche zur Verfügung und kann ab sofort beginnen.

Mit freundlichen Grüßen

\vspace{20pt}
Kartavya Niraj Dikshit

\end{document}